\section{Objetivos del trabajo}
\label{sec:objetivos-trabajo}

\subsection{Objetivo general}

El objetivo general de este Trabajo de Fin de Grado consiste en construir una
plataforma de visualizaci\'on y apoyo al an\'alisis de estrategias de trading.
Esta plataforma busca ordenar el an\'alisis del mercado a partir de datos,
reducir el tiempo dedicado a tareas repetitivas de revisi\'on, priorizar
candidatos de estudio y mantener una visi\'on
ordenada de la informaci\'on disponible. El prop\'osito no es sustituir el
criterio del usuario ni desarrollar un sistema que opere de forma totalmente
aut\'onoma. Se busca crear una base t\'ecnica que permita estudiar reglas de
entrada, contextualizarlas con datos de mercado, aprovechar mejor la evidencia
disponible y presentarla de forma clara para facilitar la toma de decisiones
supervisada.

Para ello, el trabajo busca convertir estrategias y criterios originalmente
descritos de forma discrecional en reglas que puedan calcularse, compararse y
evaluarse de manera sistem\'atica mediante distintos m\'etodos. Esta
formalizaci\'on se basa en datos intradiarios y diarios. Adem\'as, se apoya en
informaci\'on externa de contexto cuando procede, en procesos de backtesting y
en la comparaci\'on con estrategias de referencia. De este modo, la evaluaci\'on
de cada estrategia no se apoya solo en la lectura puntual de una se\~nal. Tambi\'en
se apoya en un conjunto de evidencias revisables: m\'etricas de comportamiento,
tablas comparativas, gr\'aficos, informes y registros generados durante las pruebas.

El sistema resultante se plantea como una base ampliable. Permite incorporar
nuevas estrategias al Screener, entendido como filtro de candidatos, generar
estad\'isticas sobre su comportamiento y estudiar su robustez antes de considerar
cualquier uso m\'as operativo. En este contexto, se presenta una plataforma
interactiva que agrupa gr\'aficos, Screener y an\'alisis estructural del precio
mediante WaveCount y WeaveCount. Tambi\'en integra AI Analyst, conexi\'on con MT5
y Telegram como elementos de apoyo, seguimiento y organizaci\'on de la informaci\'on.

\subsection{Objetivos espec\'ificos}

A partir del objetivo general, el trabajo se concreta en los siguientes
objetivos espec\'ificos:

\begin{enumerate}
    \setlength{\itemsep}{0.35em}
    \setlength{\parsep}{0pt}
    \item Recopilar, preparar y organizar los datos necesarios para el estudio,
    combinando series intradiarias, datos diarios e informaci\'on externa de
    apoyo cuando est\'e disponible. Este objetivo incluye el dise\~no de un flujo
    de ingesta incremental, as\'incrono y optimizado en tiempo, orientado a
    mantener actualizada la informaci\'on de mercado sin repetir procesos
    innecesarios.

    \item Formalizar estrategias de trading y criterios de contexto descritos
    originalmente de forma discrecional. Para ello, se traducen reglas de
    entrada, condiciones de mercado, niveles t\'ecnicos y filtros de revisi\'on a
    procedimientos calculables, de forma que puedan evaluarse de manera
    sistem\'atica.

    \item Dise\~nar y ejecutar pruebas hist\'oricas que permitan estudiar el
    comportamiento de las estrategias formalizadas. Estas pruebas deben
    compararse con estrategias de referencia y evitar problemas de
    \textit{look-ahead}, entendido como el uso de informaci\'on que todav\'ia no
    estar\'ia disponible en una revisi\'on real. Por ello, las se\~nales deben
    evaluarse solo con la informaci\'on que habr\'ia estado disponible en cada
    momento.

    \item Incorporar herramientas de an\'alisis estructural basadas en
    WaveCount y WeaveCount para estudiar la organizaci\'on del movimiento del
    precio. Este objetivo no busca automatizar una decisi\'on operativa, sino
    a\~nadir una capa de contexto que ayude a interpretar las se\~nales generadas
    por las estrategias.

    \item Desarrollar una plataforma interactiva de visualizaci\'on y revisi\'on de datos que permita
    consultar mercado, gr\'aficos, evidencias generadas, Screener, estado del
    sistema y relaciones entre activos desde una interfaz com\'un. La finalidad
    es reducir la dispersi\'on de informaci\'on, facilitar una revisi\'on
    supervisada de los candidatos priorizados y ayudar a identificar activos
    con alta correlaci\'on que podr\'ian concentrar riesgo de forma inadvertida,
    pese a parecer diversificados.

    \item Integrar un asistente de an\'alisis en modo solo lectura que ayude a
    preparar informes y paquetes reproducibles a partir de la informaci\'on
    disponible. Este asistente se plantea como apoyo a la revisi\'on, no como
    mecanismo de aprobaci\'on de operaciones.

    \item Verificar una capa controlada de observabilidad y ejecuci\'on demo mediante
    conexi\'on con MT5, RiskGuard y Telegram informativo. Este objetivo permite
    comprobar el flujo t\'ecnico de lectura, auditor\'ia y generaci\'on de propuestas de orden.
    Tambi\'en incluye el env\'io controlado de \'ordenes desde el bot de MT5 y el
    seguimiento informativo de eventos mediante Telegram, manteniendo fuera del
    alcance la operaci\'on aut\'onoma en producci\'on.
\end{enumerate}
